%% 
%% Technical Paper Template for CCS
%%
%% Adapted from SIGCONF Proceedings template
%%

%%
%% The first command in your LaTeX source must be the \documentclass command.
\documentclass[sigconf, pbalance]{acmart}

\usepackage{tabularx}
\usepackage{array}
\usepackage{listings}

%%
%% Inline SQL keywords. Use \sql{...}; it works inside tables and captions.
\definecolor{sqlkeyword}{HTML}{2F5D8C}
\newcommand{\sql}[1]{\texttt{\textcolor{sqlkeyword}{\bfseries #1}}}

%%
%% SQL code blocks, use \begin{lstlisting| ... \end{lstlisting}.
\lstdefinestyle{sqlstyle}{
  language=SQL,
  basicstyle=\ttfamily\footnotesize,
  keywordstyle=\color{sqlkeyword}\bfseries,
  commentstyle=\color{gray}\itshape,
  breaklines=true,
  breakatwhitespace=true,
  columns=fullflexible,
  frame=single,
  xleftmargin=0pt,
  aboveskip=6pt,
  belowskip=6pt
}
\lstset{style=sqlstyle}


%%
%% Let long \texttt identifiers break after underscores, and give TeX some
%% slack before it overfills a line.
\AtBeginDocument{%
  \let\oldunderscore\_%
  \renewcommand{\_}{\oldunderscore\allowbreak}}
\setlength{\emergencystretch}{1.5em}

\AtBeginDocument{%
  \providecommand\BibTeX{{%
    \normalfont B\kern-0.5em{\scshape i\kern-0.25em b}\kern-0.8em\TeX}}}

%%
%% Remove the ACM References statement and copyright notice.
\settopmatter{printacmref=false}
\renewcommand\footnotetextcopyrightpermission[1]{}


%%
%% Majority of ACM publications use numbered citations and references.
%% For managing citations, it is recommended to use bibliography files in BibTeX format.
%% You can then either use BibTeX with the ACM-Reference-Format style,
%% or BibLaTeX with the acmnumeric or acmauthoryear sytles, that include
%% support for advanced citation of software artefact from the
%% biblatex-software package, also separately available on CTAN.
%%
%% Look at the sample-reference.bib file enumerating your References.


%% %% %% %% %%
%% Start of the actual paper.
\begin{document}

%%
%% Paper title.
\title{A Data Warehouse and OLAP Application for the 2025 Dutch Railway Network}

%%
%% Authors
\author{Angelica Therese Clavano}
\affiliation{%
  \institution{De La Salle University}
  \city{Manila}
  \country{Philippines}}
\email{angelica\_therese\_clavano@dlsu.edu.ph}

\author{Enrique Miguel Aragon}
\affiliation{%
  \institution{De La Salle University}
  \city{Manila}
  \country{Philippines}}
\email{enrique\_aragon@dlsu.edu.ph}

\author{Jacoba Ysabelle Luna}
\affiliation{%
  \institution{De La Salle University}
  \city{Manila}
  \country{Philippines}}
\email{jacoba\_luna@dlsu.edu.ph}

\author{Zhean Robby Ganituen}
\affiliation{%
  \institution{De La Salle University}
  \city{Manila}
  \country{Philippines}}
\email{zhean\_robby\_ganituen@dlsu.edu.ph}


%%
%% By default, the full list of authors will be used in the page headers.
%% Often, this list is too long and will overlap other information printed in the page headers. 
%% This command allows the author to define a more concise list of authors' names for this purpose.
\renewcommand{\shortauthors}{Trovato et al.}


%%
%% The abstract is a short summary of the work to be presented in the article.
\begin{abstract}
This set of guidelines for formatting your research paper has been adapted from the formatting guidelines for ACM SIG Proceedings. The Abstract should consist of 150 to 250 words in a single paragraph that provides the reader with a summary of your research/project and what you intend to do. It gives a brief description of the motivation for your research/project, the problem you intend to cover, the purpose and method for conducting your study, and your results and findings.
\end{abstract}


%%
%% Keywords
%% The author(s) should pick words that accurately describe the work being presented. 
%% Separate the keywords with commas.
%% Aside from proper nouns and acronyms, only the first letter of the first keyword should be capitalized
\keywords{Data Warehouse, Dutch Railway, ETL, OLAP, Query Processing, Query Optimization}

%%
%% This command processes the author and affiliation and title information and builds the first part of the formatted document.
\maketitle

\input{chapters/introduction}
\input{chapters/data_warehouse}
\input{chapters/etl_script}
\input{chapters/olap_application}
\input{chapters/qpo}
\input{chapters/results_and_analysis}
\input{chapters/conclusion}


\bibliographystyle{ACM-Reference-Format}
\bibliography{references}

\section*{Declaration of Generative AI Usage}

% [ fr: @zrygan; @to: all
%   If you used a different model indicate in the "we used ..." and add
%   a market like (*, **, etc).
% 
%   Then mark it after \item
% ]
During the conduct of this hands-on activity, we used Claude Opus 5.5 to assist with the following tasks. The full prompts are in \texttt{dutch\_railway-docs/prompts/} of our repository.
\begin{itemize}
  \item Help understand documentation of Docker and Apache Airflow.
  \textbf{Sample prompt}: \emph{I'm a student building a data warehouse project. I'm comfortable with data modelling (star schemas, fact and dimension tables, ETL concepts) but new to [Apache Airflow / running Airflow in Docker]. Act as a teacher, not a coder: do not write code or config for me. Instead, (1) explain what [the tool / the integration] is for and what problem it solves, in a few short paragraphs; (2) map its core ideas to data modelling concepts I already know, and say where each analogy breaks down; (3) point me to the parts of the official documentation I should read first, in order; and (4) end with 3 questions I should be able to answer once I understand it, so I can check myself. If something depends on a version or setup detail, ask me instead of assuming.}
  \item \LaTeX~formatting and ACM template formatting.
  \textbf{Sample prompt}: \emph{Our report uses the ACM \texttt{acmart} class with the \texttt{sigconf} option (see \texttt{paper.tex}). Please reformat every table in \texttt{paper.tex} and \texttt{chapters/*.tex} so they follow the ACM template: use \texttt{booktabs} rules and no vertical lines; put the caption above the table, and give each table a label and a description for accessibility; make tables fit the column, using \texttt{table*} only when a table cannot fit in one column; and right-align numbers with thousands separators. Do not change any numbers or wording inside the tables. When you're done, list each table you changed and what you changed, and point out any table you weren't sure about.}
  \item Figure creation using TikZ.
  \textbf{Sample prompt (schema)}: \emph{Using the warehouse schema in our DDL (fact tables \texttt{fact\_service} and \texttt{fact\_stop}, dimension tables \texttt{dim\_date}, \texttt{dim\_time}, \texttt{dim\_station}), create a TikZ entity-relationship diagram for our ACM two-column paper. Make it a standalone file in \texttt{figures/}, so it can be compiled to PNG or PDF and included as an image. Show each table with its columns and mark primary and foreign keys. Visually separate fact tables from dimension tables, and draw one line per foreign key with 1 and N labels. Match the colours of our other figures and keep the text readable at full page width. Take the table and column names from the DDL; don't make any up. Compile the file and check that nothing overlaps.}

  \textbf{Sample prompt (query plans)}: \emph{For each query in \texttt{dutch\_railway/qpo/queries/}, create a TikZ figure for our ACM paper that compares the baseline plan with the optimized plan. Put the baseline on the left and the optimized plan on the right, as bottom-up operator trees. Highlight the operators that differ between the two plans in one accent colour, and add a one-line note under each side saying why the optimized version is faster. Use real numbers from the saved \texttt{EXPLAIN ANALYZE} output where they are available; if a plan isn't saved or looks stale, draw that side from the SQL and tell me. Use one shared style file that matches our existing figures, and make each figure a standalone file that compiles to PNG. Insert the figures into \texttt{chapters/qpo.tex} with captions and descriptions. Compile every figure, check that the layouts don't overlap, and tell me about anything you could not verify.}
\end{itemize}

Statement of Accountability. By submitting this report, we explicitly affirm the following:

\begin{itemize}
  \item We accept full accountability for the accuracy and integrity of all content in our written and oral reports. 
  \item We have examined and validated all AI-assisted contributions included in our written and oral reports.
  \item We acknowledge the inherent limitations of GenAI systems, including the potential for technical errors, factual inaccuracies, hallucinations, or biases.
  \item We fully comprehend the AI-generated outputs used in our written and oral reports. We understand that the inability to explain these concepts in our own words during evaluations will be treated as a violation of academic integrity, and we accept the consequences of any such misuse.
\end{itemize}

\section*{Record of Contribution}

% [ fr: @zrygan; @to: all
%   Task 1 & 
%   Jack's 
%   & 
%   % Ena's
%   & 
%   % Koby's
%   & 
%   % Zen's
%   \\
% ]
\begin{table}[htbp]
\centering
\scriptsize
\caption{Record of member contributions}
\label{tab:contributions}
\Description{Record of individual group member contributions across project tasks.}
\renewcommand{\arraystretch}{1.2}
\begin{tabularx}{\columnwidth}{>{\raggedright\arraybackslash}p{2.0cm} *{4}{>{\centering\arraybackslash}X}}
\toprule
\textbf{Task} & \textbf{Clavano} & \textbf{Aragon} & \textbf{Luna} & \textbf{Ganituen} \\
\midrule
\textbf{Data Warehouse} & 
Database schema/design & 
Database design & 
Database design & 
Exploratory data analysis \\
\addlinespace[4pt]
\textbf{ETL Pipeline} & 
--- & 
--- & 
--- & 
Pipeline implementation \\
\addlinespace[4pt]
\textbf{OLAP Application} & 
App development and deployment & 
--- & 
Query ideation and baseline implementation & 
--- \\
\addlinespace[4pt]
\textbf{Functional Testing} & 
--- & 
Performed functional testing & 
--- & 
Test implementation \\
\addlinespace[4pt]
\textbf{Performance Testing} & 
--- & 
Optimized queries for reports 1, 2, 4  & 
Optimized query for report 3 & 
Test implementation \\
\addlinespace[4pt]
\textbf{Report} & 
Wrote Application Development section in \ref{sec:olap_application} & 
Wrote \ref{sec:data_warehouse} and \ref{sec:qpo} &
Wrote sections 4.1 through 4.4 in \ref{sec:olap_application}  & 
Wrote \ref{sec:introduction} and \ref{sec:etl_script} \\
\bottomrule
\end{tabularx}
\end{table}
\clearpage
\appendix
\input{chapters/appendix}

\end{document}
